\section{Related Work}
\label{sec:related-work}

This section situates the proposed work within four strands of prior research: agricultural pest and disease diagnosis expert systems, the intersection of large language models and knowledge graphs in agriculture, Model Context Protocol and configuration-driven architectures, and Rust-native agent frameworks. Each subsection concludes with a critical summary that identifies the specific limitation our study aims to overcome.

\subsection{Agricultural Pest and Disease Diagnosis Expert Systems}

Traditional approaches to computer-aided crop protection have relied heavily on explicit, hand-curated knowledge representations. Jearanaiwongkul et al. \cite{jearanaiwongkul2021riceman} developed RiceMan, an ontology-based expert system that encodes rice-disease knowledge through the Rice Disease Ontology (RiceDO) and performs symptom-to-disease mapping via Semantic Web Rule Language (SWRL) reasoning. The system offers transparent, interpretable diagnostic trails that satisfy the explainability requirements of agricultural extension services; yet its knowledge base is tightly coupled to the rice domain, making cross-crop adaptation prohibitively expensive because every new crop would require the construction of a new ontology and rule set. Similarly, Alharbi et al. \cite{alharbi2024agrodss} proposed AgrODSS, which combines a plant-disease ontology (PDP-O) with an evidence-based explanation model (EBEM). Although expert evaluation reached an accuracy of 80.66\% on held-out diagnostic cases, the underlying SPARQL queries and OWL axioms would require substantial restructuring for other cultivation systems or geographic regions with different pest pressures.

More recent efforts have sought to harness the flexible reasoning capabilities of LLMs while retaining the structured fidelity of agricultural knowledge bases. The AgenticGraphRAG framework \cite{agenticgraphrag2025} integrates ReAct-style agent orchestration with dual retrieval---dense vector search plus Cypher graph queries---to achieve 86--89\% accuracy on rice pest and disease diagnosis. Nevertheless, its system prompts, retrieval module templates, and tool chains remain hard-coded around rice-specific schemas and symptom vocabularies. At the frontier of diagnostic accuracy, Chat Demeter \cite{chatdemeter2025} deploys a CNN--Transformer vision backbone within a four-agent collaborative pipeline encompassing planning, reasoning, evaluation, and visualization agents, reporting an impressive 99.50\% classification accuracy on plant-leaf images. Yet this performance comes at the price of extreme vertical customization: neither the vision model nor the multi-agent workflow can be redeployed to a new crop without domain-specific re-engineering of both the neural architecture and the inter-agent communication protocol. Complementary work by Chelloug et al. \cite{chelloug2023mdtwlstm} and Rodr{\'i}guez-Garc{\'i}a et al. \cite{rodriguez2021kbs} further demonstrates that early knowledge-based systems, while foundational, share the same limitation of crop-specific knowledge engineering.

\textbf{Critical summary.} Despite achieving respectable and sometimes state-of-the-art accuracy, existing diagnosis systems---whether ontology-driven, LLM-augmented, or multi-agent vision pipelines---are predominantly vertical applications optimized for a single crop or diagnostic task. What is conspicuously missing from the literature is a horizontal agent platform in which agricultural domain logic, diagnostic schemas, and tool inventories can be injected declaratively at runtime, without rewriting rules, retraining models, or refactoring agent workflows.

\subsection{Agricultural Large Language Models and Knowledge Graphs}

The fusion of knowledge graphs (KGs) and LLMs has attracted growing attention in the agricultural informatics community as a means of mitigating hallucination, improving factual grounding, and enabling traceable question answering. Gong and Li \cite{gong2025kgllm} provide a comprehensive survey of KG--LLM complementarity in agriculture, arguing that structured knowledge and parametric language models can be synergistically combined to overcome both the high manual cost of KG construction and the factual unreliability of standalone LLMs. Empirical support for this thesis comes from Zhao et al. \cite{zhao2024llmkgagriculture}, who demonstrate that injecting KG-derived facts into LLM prompts significantly improves both factual accuracy and domain relevance in plant-disease detection tasks compared with ungrounded baselines.

On the model-development side, the KGLLM system \cite{kgllm2024} employs information-entropy filtering and explicit constraint decoding guided by a knowledge graph, achieving an average BertScore improvement of 9.84\% over vanilla LLM generation while markedly reducing the incidence of hallucinated pesticide recommendations. AgriGPT \cite{yang2025agrigpt} represents the first open-source ecosystem dedicated to agricultural LLMs, releasing the Agri-342K instruction-following dataset, a Tri-RAG retrieval pipeline, and the AgriBench-13K evaluation benchmark. These resources have lowered the barrier to entry for agricultural NLP research, yet the reference implementation assumes cloud-hosted inference. Meanwhile, AgroLLM \cite{samuel2025agrollm} introduces domain-constrained retrieval-augmented generation through a Domain Knowledge Prompting Library (DKPL), automatically extracting semantic vocabularies, causal rules, and constraint libraries from agricultural textbooks to reach 95.2\% accuracy on farmer advisory tasks. Rezayi et al. \cite{rezayi2022agribert} provided an early foundation for this line of work by showing that continued pre-training on agricultural corpora yields substantial gains in domain text understanding.

\textbf{Critical summary.} While these works convincingly show that KG grounding can reduce hallucination and boost answer accuracy, the majority are engineered for cloud deployment and assume continuous access to large GPU servers or proprietary LLM APIs. Research on local-first adaptation---running the full KG--LLM pipeline on commodity edge devices, offline workstations, or low-power rural computers---remains conspicuously sparse in the agricultural literature, creating a significant barrier for practitioners in bandwidth-constrained environments.

\subsection{Model Context Protocol and Configuration-Driven Architectures}

The Model Context Protocol (MCP), introduced and maintained by Anthropic \cite{anthropic2024mcp}, defines a standardized JSON-RPC~2.0 interface through which LLM-based applications can discover, invoke, and compose external tools. By separating the client (the reasoning application) from the server (the tool provider), MCP promises to reduce the integration burden that has historically plagued agent deployments. In the industrial Internet-of-Things (IoT) domain, the IoT-MCP framework \cite{iotmcp2025} proposes a three-layer architecture---comprising a Local Host, a Datapool and Connection Server, and distributed IoT Devices---that translates natural-language commands into standardized JSON control messages for edge sensors and actuators. Complementing this vision, production-grade MCP servers for manufacturing environments now expose MQTT and Modbus interfaces as unified AI-orchestrable APIs \cite{iotedgemcp2025}, enabling predictive maintenance and anomaly detection workflows without bespoke API integration for each device vendor.

In agriculture, declarative configuration has been explored primarily at the infrastructure and cyber-physical systems level. The AGRARIAN project \cite{agrarian2025} leverages lightweight Kubernetes (K3s) and GitOps principles to orchestrate edge nodes and AI containers in cyber-physical farming systems, achieving dynamic service deployment through declarative manifests rather than imperative shell scripts. Groeneveld et al. \cite{groeneveld2021dsl} developed a domain-specific language (DSL) framework for farm-management information systems (FMIS), explicitly separating functionality, sector, and configuration concerns to reduce communication overhead between farmers and software engineers. These precedents underscore the practical value of declarative abstraction in complex agro-ecological settings. The local-first software movement \cite{kleppmann2021localfirst} further provides a philosophical and technical foundation for edge-native applications that prioritize data sovereignty and offline availability over cloud dependence.

\textbf{Critical summary.} Although MCP and declarative configuration have shown substantial promise in industrial IoT, manufacturing automation, and infrastructure orchestration, no standardized tool-registration schema, context-exchange format, or offline transport profile currently exists for agricultural vertical domains. The gap between general-purpose protocol specifications and the specific requirements of farm-level knowledge management---such as pest-ontology bindings, crop-specific JSON schemas, and agronomic context serialization---therefore remains unbridged.

\subsection{Rust-Native Agent Frameworks}

The systems-programming language Rust has garnered increasing interest in the agent-framework community owing to its memory-safety guarantees, deterministic resource management, and zero-cost abstractions---all of which are particularly attractive for long-running edge services. Rig \cite{playgrounds2025rig}, developed by Playgrounds, emphasizes scalability, WebAssembly portability, and first-class vector-store integration, yet it provides neither a native terminal user interface (TUI) nor built-in subagent hierarchies, limiting its utility for interactive, multi-step agricultural advisory sessions. RusticAI \cite{rusticai2025}, inspired by PydanticAI, focuses on real-time voice interaction and deferred tool execution; while it supports MCP toolsets, it lacks a subagent permission-inheritance model that would be necessary for complex, hierarchical task decomposition. ADK-Rust \cite{zavora2025adkrust} pursues a model-agnostic, deployment-independent architecture for voice agents, but again omits both a TUI and hierarchical subagent capabilities, positioning it more as a backend inference orchestrator than an end-user knowledge-management workstation.

The broader MCP ecosystem, catalogued in community-curated lists \cite{awesomemcp2025}, offers reference implementations for filesystem access, memory stores, and web-fetch tools; however, no agriculture-specific MCP server or client reference implementation is currently listed, reflecting the nascent state of the field.

\textbf{Critical summary.} To the best of our knowledge, no existing Rust agent platform simultaneously offers a native TUI, hierarchical subagent support, and a standards-compliant MCP client within a local-first deployment model. This absence motivates the development of Clarity as the application-layer kernel in the architecture proposed by the present study, and it underscores the novelty of integrating such a kernel with a local agricultural knowledge base via MCP.
